problem:  
Let \( M^n \) be a compact, simply connected Riemannian manifold with positive sectional curvature. Suppose an effective, isometric torus \( T^k \)-action operates on \( M \).  
Assume that for every circle subgroup \( S^1 \subset T^k \), there exists at least one connected component \( F_{S^1} \subset M^{S^1} \) of its fixed-point set whose codimension satisfies  
\[
\mathrm{codim}(F_{S^1}) \le c,
\]  
for some constant \( c \) independent of the subgroup.  

**Conjecture / Problem:**  
Determine whether there exists an explicit linear (or sublinear) function \( f: \mathbb{N} \to \mathbb{N} \) such that if \( c \le f(n) \), then \( M \) must be diffeomorphic to a compact rank-one symmetric space (CROSS).  
In particular, is it true that \( f(n) = n/6 \) (or any comparable bound) suffices?  
Furthermore, under the same assumption, does the existence of small-codimension fixed-point sets for all \( S^1 \subset T^k \) imply periodicity in the rational cohomology ring of \( M \)?

Why is it a "good" problem:  
This problem directly intertwines *quantitative geometric symmetry* with *topological rigidity*, extending the known symmetry rank rigidity program in a genuinely new direction. Unlike the classical Grove–Searle and Wilking theorems, which focus on **global rank extremality** or the **existence** of a single small-codimension totally geodesic submanifold, the current problem imposes a *uniform codimension constraint across all circle subgroups*—a much stronger global geometric coherence condition.  

It would force an exploration of how widespread geometric “low-codimension symmetry” interacts with the topology of positively curved manifolds, potentially leading to new periodicity or vanishing results beyond Wilking’s connectedness lemma and the Kennard periodicity framework. Determining any explicit bound \( f(n) \) would refine our understanding of how symmetry distributes curvature and shape on manifolds of positive curvature.  

The conjecture is clear, self-contained, and lies at the intersection of Riemannian geometry, transformation group theory, and global topology—an ideal research-level problem: simple in its clauses yet deeply intertwined with multiple major geometric rigidity mechanisms.